\documentclass[11pt]{article}

\usepackage[a4paper,margin=1in]{geometry}
\usepackage{amsmath,amssymb,amsthm,mathtools}
\usepackage{booktabs}
\usepackage[hidelinks]{hyperref}

\newtheorem{theorem}{Theorem}[section]
\newtheorem{proposition}[theorem]{Proposition}
\newtheorem{lemma}[theorem]{Lemma}
\newtheorem{corollary}[theorem]{Corollary}
\theoremstyle{definition}
\newtheorem{remark}[theorem]{Remark}

\newcommand{\R}{\mathbb{R}}
\newcommand{\Fr}{\mathcal{A}}
\newcommand{\Om}{\Omega}
\newcommand{\Otilde}{\widetilde{\Omega}}
\newcommand{\eps}{\varepsilon}

\title{Sharp Quantitative Faber--Krahn Stability\\
       via a Contradiction-Free Route}
\author{Plan~1 / Route~A: assembled document}
\date{}

\begin{document}
\maketitle

\begin{abstract}
We give a proof of the sharp quantitative Faber--Krahn inequality
\[
   \Bigl(\frac{|\Om|}{|B_1|}\Bigr)^{2/N}
   \bigl(\lambda_1(\Om)-\lambda_1(B^*)\bigr)
   \ge c_{\rm FK}(N)\,\Fr(\Om)^2,
\]
with sharp exponent $2$ on the Fraenkel asymmetry $\Fr$ (equivalently
$1/2$ on the scale-normalised deficit), in which the stability constant
$c_{\rm FK}(N)$ is an explicit composition of named constants. The
sequential-selection / compactness contradiction step in the proof of
Brasco--De~Philippis--Velichkov \cite{BDV} is replaced by a
single-set quantitative selection principle, and every subsequent step
is constructive. The proof is organised into seven self-contained
building blocks, each delivered as a separate \LaTeX{} document:
\begin{itemize}
   \item \texttt{SelectionPrinciple.tex} --- single-set quantitative
   selection;
   \item \texttt{GraphEntry.tex} --- BDV asymmetry $\to$ global
   $C^{1,\gamma}$ normal graph;
   \item \texttt{SchauderInterpolation.tex} --- $C^{1,\gamma}$ $\to$
   small $C^{2,\gamma_0}$;
   \item \texttt{NearlySphericalClosure.tex} --- BDV Theorem~3.3
   delivers $Q_\alpha\ge c_*(N)$;
   \item \texttt{SaintVenantAssembly.tex} --- small + far regimes glue
   to sharp Saint--Venant stability inside a fixed ball $B_R$;
   \item \texttt{BoundedReduction.tex} --- removes the
   $R$-dependence by reducing the global case to the bounded one via
   the constructive BDV truncation (\cite[\S5]{BDV});
   \item \texttt{KohlerJobinTransfer.tex} --- Saint--Venant
   $\to$ Faber--Krahn via Kohler--Jobin.
\end{itemize}
The introduction sets the result in the literature
\cite{FMP,BDV,Brasco2014}; this master document gives the chain of
labelled theorems and the final assembled statement.
\end{abstract}

\tableofcontents

\section{Introduction}\label{sec:intro}

\subsection{The Faber--Krahn inequality}\label{ssec:FK}

Let $\Om\subset\R^N$, $N\ge2$, be an open set of finite Lebesgue measure
and let
\[
   \lambda_1(\Om)\;=\;\min_{u\in H^1_0(\Om)\setminus\{0\}}
                    \frac{\int_\Om|\nabla u|^2\,dx}{\int_\Om u^2\,dx}
\]
be the first Dirichlet eigenvalue of the Laplacian on $\Om$. The
\emph{Faber--Krahn inequality} asserts that among open sets of a
prescribed volume the ball uniquely minimises $\lambda_1$: writing
$B^*$ for the ball with $|B^*|=|\Om|$,
\begin{equation}\label{eq:FKsharp}
   \lambda_1(\Om)\;\ge\;\lambda_1(B^*),
\end{equation}
with equality if and only if $\Om$ is a ball (up to a set of capacity
zero). This was conjectured by Lord Rayleigh in the late nineteenth
century in the planar case and proved in general dimension by
Faber~\cite{Faber1923} and Krahn~\cite{Krahn1925} using Schwarz
symmetrisation and the P\'olya--Szeg\H{o} principle. Rigidity reflects
the fact that the only equality cases of Schwarz symmetrisation in the
gradient norm are radially symmetric functions; modern accounts
appear in~\cite{KawohlBook,HenrotBook,BucurButtazzo}.

Once~\eqref{eq:FKsharp} has been proved with rigidity, it is natural
to ask for a \emph{stability} statement: if $\lambda_1(\Om)$ is only
slightly larger than $\lambda_1(B^*)$, then $\Om$ should be in a
quantitative sense close to a ball. The natural way to measure the
deviation from spherical symmetry is the \emph{Fraenkel asymmetry}
\begin{equation}\label{eq:Fraenkel}
   \Fr(\Om)\;=\;\inf_{x_0\in\R^N}
   \frac{|\Om\,\Delta\,(B^*+x_0)|}{|B^*|}\;\in\;[0,2),
\end{equation}
which vanishes precisely on translates of the ball. A quantitative
Faber--Krahn inequality is then an estimate of the form
\begin{equation}\label{eq:qFKabstract}
   \Bigl(\frac{|\Om|}{|B_1|}\Bigr)^{2/N}
   \bigl(\lambda_1(\Om)-\lambda_1(B^*)\bigr)
   \;\ge\;c\,\omega(\Fr(\Om)),
\end{equation}
for some modulus $\omega:[0,2)\to[0,\infty)$ vanishing only at~$0$.
The first systematic study of such bounds is due to Hansen and
Nadirashvili~\cite{HansenNadirashvili1994}, with parallel
contributions by Melas~\cite{Melas1992}; both gave
bounds of the shape~\eqref{eq:qFKabstract} only for restricted
classes of sets and with a non-sharp modulus. The first stability
result valid for every open set in terms of the Fraenkel asymmetry was
obtained by Bhattacharya~\cite{Bhattacharya2001} in dimension $N=2$
(exponent~$3$). Fusco, Maggi and Pratelli~\cite{FMPLaplace}, exploiting
their sharp quantitative isoperimetric inequality~\cite{FMP}, extended
this to all dimensions with exponent~$4$ on~$\Fr$; see
also~\cite{Brasco2014,BraDeP2017}.

\subsection{The sharp exponent}\label{ssec:sharpexp}

The optimal modulus in~\eqref{eq:qFKabstract} is dictated by the
family of linear ellipsoidal perturbations of the ball,
$\Om_\eps=\{(x',x_N)\in\R^N:|x'|^2+(1+\eps)x_N^2\le1\}$,
$0<\eps\ll1$, for which a direct computation yields
\[
   \Fr(\Om_\eps)\;\simeq\;\eps,\qquad
   \lambda_1(\Om_\eps)-\lambda_1(B^*)\;\simeq\;\eps^2;
\]
see~\cite[\S1.2]{BDV}. Consequently the sharp exponent is~$2$ on the
asymmetry, equivalently $1/2$ on the spectral deficit, and the sharp
quantitative Faber--Krahn inequality is the conjecture
\begin{equation}\label{eq:sharpconj}
   \Bigl(\frac{|\Om|}{|B_1|}\Bigr)^{2/N}
   \bigl(\lambda_1(\Om)-\lambda_1(B^*)\bigr)
   \;\ge\;c(N)\,\Fr(\Om)^2.
\end{equation}
The exponent~$2$ is natural because the second variation of $\lambda_1$
around the ball is non-degenerate on the subspace of volume- and
barycentre-preserving normal perturbations.
The conjecture~\eqref{eq:sharpconj} sits at the spectral end of a
family of sharp quantitative problems (isoperimetric, Sobolev,
log-Sobolev, P\'olya--Szeg\H o, Wulff, Riesz potentials), the
prototype being the sharp quantitative isoperimetric inequality of
Fusco--Maggi--Pratelli~\cite{FMP}, with later proofs and refinements
by Cicalese--Leonardi~\cite{CL2012},
Figalli--Maggi--Pratelli~\cite{FigalliMaggiPratelli}, and
Fusco--Julin~\cite{FuscoJulin}; see~\cite{FuscoSurvey,Maggi2008} for
surveys.

\subsection{The Brasco--De~Philippis--Velichkov theorem}\label{ssec:BDV}

The conjecture~\eqref{eq:sharpconj} was settled in 2015 by Brasco,
De~Philippis and Velichkov~\cite{BDV}:

\begin{theorem}[Brasco--De~Philippis--Velichkov; {\cite[Main Theorem]{BDV}}]
\label{thm:BDV}
For every $N\ge2$ there is a constant $\sigma=\sigma(N)>0$ such that
\[
   |\Om|^{2/N}\lambda_1(\Om)-|B^*|^{2/N}\lambda_1(B^*)
   \;\ge\;\sigma\,\Fr(\Om)^2
\]
for every open set $\Om\subset\R^N$ of finite measure.
\end{theorem}

The strategy of~\cite{BDV} has two main pillars. First, an application
of the Kohler--Jobin rearrangement inequality~\cite{KohlerJobin,Brasco2014}
reduces~\eqref{eq:sharpconj} at fixed volume to the analogous
\emph{sharp Saint--Venant stability}
\begin{equation}\label{eq:sharpSV}
   E(\Om)-E(B_1)\;\ge\;\sigma\,\Fr(\Om)^2,\qquad |\Om|=|B_1|,
\end{equation}
where $E(\Om)=\min\{\tfrac12\int|\nabla u|^2-\int u:u\in H^1_0(\Om)\}=
-\tfrac12 T(\Om)$. Second, \eqref{eq:sharpSV} is proved by the
\emph{selection principle} introduced by Cicalese and
Leonardi~\cite{CL2012} for the isoperimetric inequality and refined by
Acerbi--Fusco--Morini~\cite{AFM} in the nonlocal isoperimetric
setting. Arguing by contradiction, one supposes a sequence $\Om_j$
violates~\eqref{eq:sharpSV} with $\Fr(\Om_j)\to0$ and constructs an
improved sequence $U_j$ by penalising the energy along BDV's smooth
asymmetry $\alpha$:
\begin{equation}\label{eq:BDVpenalised}
   U_j\;\in\;\arg\min\Bigl\{\,E(\Om)
      +\sqrt{\eps_j^2+\sigma(\alpha(\Om)-\eps_j)^2}\,:\,|\Om|=|B_1|\Bigr\}.
\end{equation}
The Alt--Caffarelli free-boundary regularity machinery~\cite{AC1981}
applied to the perturbed Bernoulli problem~\eqref{eq:BDVpenalised}
upgrades $L^1$ convergence $U_j\to B_1$ to $C^{2,\gamma}$ convergence
of $\partial U_j$ as normal graphs. Once $U_j$ is a smooth normal
graph over $B_1$ with small $C^{2,\gamma}$ norm, the explicit
second-order Fuglede-type expansion~\cite{Fuglede,BDV} of the
Saint--Venant energy --- BDV's nearly spherical Theorem~3.3 --- yields
the desired contradiction.

\subsection{The cost of the contradiction step}\label{ssec:cost}

The proof of Theorem~\ref{thm:BDV} just sketched is conceptually clean
but inherently \emph{non-constructive}: the constant $\sigma(N)$ in
the main theorem is produced by a compactness--contradiction argument
applied to the penalised sequence $U_j$. Three layers of
non-effectiveness compound:
\begin{enumerate}
\item[\rm(a)] The selection step~\eqref{eq:BDVpenalised} works only
``for $j$ large'' and the threshold is not localised.
\item[\rm(b)] The smooth convergence $\partial U_j\to\partial B_1$ in
$C^{2,\gamma}$ is obtained as a subsequential limit; no explicit
threshold is given at which the graph entry holds.
\item[\rm(c)] The contradiction itself uses the limit set $U_\infty$,
not any specific $U_j$, so even if (a) and (b) were quantified the
final constant $\sigma(N)$ would still be inaccessible.
\end{enumerate}
As a consequence, although Theorem~\ref{thm:BDV} is fully proved as a
mathematical statement, the constant $\sigma(N,R)$ (restricted to sets
in a fixed ball $B_R$) is not extractable from the proof, even in
principle.

This phenomenon is well known in the literature on sharp quantitative
stability and is not specific to~\cite{BDV}: every proof based on
penalisation + selection + Alt--Caffarelli regularity that we are
aware of shares the same feature; see the
surveys~\cite{FuscoSurvey,Maggi2008}.

\subsection{Main result and contribution}\label{ssec:contribution}

We give a proof of the sharp quantitative Faber--Krahn inequality in
which the contradiction step is replaced by a single-set quantitative
selection map and every subsequent step is constructive. Our main
theorem is the following.

\begin{theorem}[Sharp quantitative Faber--Krahn, constructive form]
\label{thm:mainintro}
Let $N\ge2$. There exists a constant $c_{\rm FK}(N)>0$, depending
only on the dimension, given explicitly by the closed-form
expression~\eqref{eq:cFKdef} below as a composition of the named
constants of Theorems~\ref{thm:S}--\ref{thm:KJ}, such that
\[
   \Bigl(\frac{|\Om|}{|B_1|}\Bigr)^{2/N}
   \bigl(\lambda_1(\Om)-\lambda_1(B^*)\bigr)
   \;\ge\;c_{\rm FK}(N)\,\Fr(\Om)^2
\]
for every open $\Om\subset\R^{N}$ of finite positive measure. The exponent~$2$
on $\Fr$ is sharp.
\end{theorem}

The novelty is not the inequality itself --- which is~\cite{BDV} ---
but its \emph{constant}. Theorem~\ref{thm:mainintro} provides a
\emph{truly quantitative} version of the BDV theorem:
$c_{\rm FK}(N)$ is presented as an explicit composition of
$N$-dependent quantities, each of which is the named constant of a
quantitative ingredient (BDV Lemmas~4.2, 4.5, 4.9, 4.12, 4.16, 5.1,
5.2, 5.3, BDV Theorems~3.3 and~4.18; the dimensional constants
$L_N,T_N$; the Lieberman--Schauder constant; the
Gagliardo--Nirenberg interpolation constant). We emphasise that we do
not compute $c_{\rm FK}(N)$ to a number: the BDV black boxes
themselves are not numerically evaluated. What we deliver is a
contradiction-free chain whose constants could in principle be
computed if the BDV constants were.

The conceptual gain is the following. In~\cite{BDV} the BDV penalised
functional~\eqref{eq:BDVpenalised} is run along a sequence of
near-extremals to produce a limit set. Here we run it \emph{once}, on
a single near-extremal $\Om_0$, with the penalisation tuned to
$\Om_0$ itself:
\[
   J_{\Om_0}(U)\;:=\;
   E(U)+f_{\hat\eta}(|U|)+\sqrt{\eps^2+\tau^2(\alpha(U)-\eps)^2},
   \qquad \eps:=\alpha(\Om_0),\ \tau=q^{1/4},
\]
where $q:=Q_\alpha(\Om_0)=(E(\Om_0)-E(B_1))/\alpha(\Om_0)$ is the
Saint--Venant quotient of $\Om_0$. The direct method produces a single
minimiser $U_*$ whose volume-normalised companion $\Otilde$ retains
the same asymmetry order as $\Om_0$ and whose quotient is at most
twice that of $\Om_0$ up to an $(N,R)$-explicit factor. The BDV
regularity package, applied to this one set, then provides a
$C^{1,\gamma}$ normal graph parametrisation of $\partial\Otilde$ at
an explicit threshold; the Kinderlehrer--Nirenberg hodograph
straightening~\cite{KN} with Lieberman's oblique Schauder
bootstrap~\cite{Li86} promote this to a uniform $C^{5,\gamma}$ bound;
Gagliardo--Nirenberg interpolation between $\|g\|_{L^\infty}$ (small)
and $\|g\|_{C^{5,\gamma}}$ (large but controlled) reaches the small
$C^{2,\gamma_0}$ regime in which BDV's nearly spherical Theorem~3.3
applies directly to $\Otilde$. No subsequence, no limit, no
contradiction.

\subsection{Structure of the paper}\label{ssec:structure}

The argument is organised as a chain of seven self-contained building
blocks, each proved in a companion file and quoted as a labelled
theorem in Section~\ref{sec:chain}. Schematically,
\[
   \underbrace{\text{(S)}}_{\text{selection}}\;\to\;
   \underbrace{\text{(G)}}_{\text{graph entry}}\;\to\;
   \underbrace{\text{(I)}}_{\text{Schauder + interpolation}}\;\to\;
   \underbrace{\text{(N)}}_{\text{nearly spherical}}
\]
operates in the small-asymmetry regime, producing the linear bound
$E(\Om_0)-E(B_1)\ge q_*(N,R)\alpha(\Om_0)$ on
$\{\alpha\le\alpha_{\rm small}\}$. Combined with BDV Lemma~4.2 this
yields the small-asymmetry portion of sharp Saint--Venant stability.
The far regime is covered by BDV's rearrangement-based Lemma~5.2 (the
suboptimal $\Fr^4$ bound, which contains no contradiction); the two
regimes are glued at an explicit threshold supplied by BDV
Lemma~4.2. The \textsc{SaintVenantAssembly} block
(Theorem~\ref{thm:SV}) gives~\eqref{eq:sharpSV} with $c_{\rm SV}(N,R)$
explicit. The Kohler--Jobin rearrangement transfer
(Theorem~\ref{thm:KJ}), in the constructive form due to
Brasco~\cite{Brasco2014}, takes Saint--Venant stability to
Faber--Krahn stability with a multiplier $4L_N/((N+2)T_N)$ that is
universal in $N$. Section~\ref{sec:assembled} states and proves the
assembled Theorem~\ref{thm:final}, and Section~\ref{sec:replaces}
summarises what replaces the contradiction.

\subsection{Related literature}\label{ssec:related}

Sharp quantitative stability of geometric and spectral inequalities
is a vast subject. The prototype is the sharp quantitative
isoperimetric inequality, first proved by
Fusco--Maggi--Pratelli~\cite{FMP} by symmetrisation, then reproved by
Cicalese--Leonardi~\cite{CL2012} via the selection principle, and by
Figalli--Maggi--Pratelli~\cite{FigalliMaggiPratelli} via optimal
transport; Fusco--Julin~\cite{FuscoJulin} gave a strong form involving
an oscillation index. Sharp quantitative
P\'olya--Szeg\H{o}, Sobolev and log-Sobolev inequalities were
established by Cianchi--Fusco--Maggi--Pratelli~\cite{CFMP}. Stability
for Wulff shapes is due to Figalli--Maggi--Pratelli and
Neumayer~\cite{NeumayerWulff}. In the spectral direction,
Brasco--Pratelli~\cite{BrascoPratelli} treated the second eigenvalue
and the Szeg\H{o}--Weinberger inequality; recent work of
Allen--Kriventsov--Neumayer~\cite{AKN} addresses Faber--Krahn-type
inequalities via Alt--Caffarelli--Friedman monotonicity formulas.

Two strands of recent work directly motivate the present paper.
First, the use of regularity theory \emph{quantitatively}: the
realisation that free-boundary regularity packages underlying
selection-principle proofs can often be made effective on a single
set if one is content with $(N,R)$-explicit constants. Second, the
technique of Fuglede-style expansion for energy functionals around
the ball~\cite{Fuglede,AFM,BDV}, which converts nearly spherical
stability into a coercivity statement for a self-adjoint operator on
$\partial B_1$. The present paper combines these ideas with the
specific BDV proof: the BDV penalised functional is run once, the BDV
regularity package is applied as a black box, the BDV nearly
spherical theorem is invoked directly on the single selected set,
and the resulting Saint--Venant stability is transferred to
Faber--Krahn via the Brasco--Kohler--Jobin rearrangement. The outcome
is a constructive constant chain for the sharp Faber--Krahn
inequality.

\begin{remark}[Reduction to bounded sets]
The chain (S)$\to$(G)$\to$(I)$\to$(N) of
Theorems~\ref{thm:S}--\ref{thm:N} produces the sharp Saint--Venant
stability inside a fixed ball $B_R$ with constant $c_{\rm SV}(N,R)$.
The \textsc{BoundedReduction} block (Theorem~\ref{thm:BR}) reduces
the global case to the bounded case at the dimensional radius
$R_{*}(N):=\max\{d(N),2\}$, using BDV's constructive truncation
lemma~\cite[Lemma~5.3]{BDV}. The constants $d(N),\delta(N),C_{9}(N)$
of that lemma are themselves obtained by De~Giorgi-type iteration and
a cut-off competitor, with no contradiction step. As a consequence
the constant $c_{\rm SV}^{\rm glob}(N)$ (and hence $c_{\rm FK}(N)$)
depends only on $N$.
\end{remark}

\section{Notation and standing hypotheses}\label{sec:notation}

Fix $N\ge2$ and $R\ge2$. We use the standard normalisations:
\begin{itemize}
\item $u_\Om\in H^1_0(\Om)$ is the torsion potential, $-\Delta u_\Om=1$
in $\Om$, $u_\Om=0$ on $\partial\Om$;
\item $T(\Om)=\int_\Om u_\Om\,dx$, the torsional rigidity;
\item $E(\Om)=-T(\Om)/2$, the BDV-normalised Saint--Venant energy:
\[
E(\Om)=\min_{v\in W^{1,2}_0(\Om)}\Bigl(\tfrac12\int|\nabla v|^2-\int v\Bigr);
\]
\item $\lambda_1(\Om)$, the first Dirichlet eigenvalue of $-\Delta$;
\item $\Fr(\Om)$, the Fraenkel asymmetry,
$\inf_{x_0}|\Om\Delta(B^*+x_0)|/|B^*|$, where $B^*$ is the ball with
$|B^*|=|\Om|$; $\Fr\in[0,2)$;
\item $\alpha(\Om)$, the BDV smooth asymmetry of
\cite[Def.~4.1]{BDV};
\item $Q_\alpha(\Om)=(E(\Om)-E(B_1))/\alpha(\Om)$ at fixed volume
$|\Om|=|B_1|=\omega_N$.
\end{itemize}
$L_N=\lambda_1(B_1)$ and $T_N=T(B_1)=|B_1|/(N(N+2))$ are universal
dimensional constants. We write $B_R$ for the open ball of radius $R$
about $0$.

\section{The chain of building blocks}\label{sec:chain}

Each block is stated here in the form it is proved in the
corresponding companion file. Constants depend only on $(N,R)$ unless
indicated.

\subsection{Block I: single-set quantitative selection}

\begin{theorem}[\textsc{SelectionPrinciple}; see
\texttt{SelectionPrinciple.tex}]\label{thm:S}
There exist $\eps_{\rm sel}(N,R), q_{\rm sel}(N,R)>0$ and
$C_{\rm sel}(N,R)>0$ such that for every quasi-open
$\Om_0\subset B_R$ with $|\Om_0|=|B_1|$,
$\eps:=\alpha(\Om_0)\le\eps_{\rm sel}$, and
$q:=Q_\alpha(\Om_0)\le q_{\rm sel}$, the penalised functional
$J(U)=F_{\hat\eta}(U)+\sqrt{\eps^2+\tau^2(\alpha(U)-\eps)^2}$
with $\tau=q^{1/4}$ admits a minimiser $U_*\subset B_R$ whose
volume-normalised companion
$\Otilde:=r_*^{-1}U_*$, $r_*=(|U_*|/|B_1|)^{1/N}$, satisfies:
\begin{enumerate}
\item[\rm(i)] $|\Otilde|=|B_1|$, $\tfrac12\eps\le\alpha(\Otilde)\le\tfrac32\eps$;
\item[\rm(ii)] $Q_\alpha(\Otilde)\le 2C_{\rm sel}(N,R)\,Q_\alpha(\Om_0)$;
\item[\rm(iii)] $\Otilde$ inherits the BDV penalised-minimiser regularity
package (density estimates, Lipschitz nondegeneracy, smooth Bernoulli
coefficient);
\item[\rm(iv)] the Bernoulli law on $\partial\Otilde$ takes the
volume-rescaled form $|\nabla\widetilde u|^2=\widetilde\Lambda+\widetilde
a_\sigma\widetilde\Psi$.
\end{enumerate}
The proof is constructive: existence is by the direct method (BDV
Lemma~4.6 used as a quantitative black box), and parts~(i)--(iv) are
explicit estimates from the BDV regularity package and the volume
scaling $E(\mu V)=\mu^{N+2}E(V)$.
\end{theorem}

\subsection{Block II: graph entry}

\begin{theorem}[\textsc{GraphEntry}; see \texttt{GraphEntry.tex}]\label{thm:G}
There exist
$\alpha_{\rm graph}(N,R)$, $M_1(N,R)$, $C_H''(N,R)>0$ and
$\gamma=\gamma(N,R)\in(0,1)$ such that every selected
$\Otilde\subset B_R$ (Theorem~\ref{thm:S}) with
$\alpha(\Otilde)\le\alpha_{\rm graph}$ is, after translation by its
Fraenkel-optimal centre, a global $C^{1,\gamma}$ normal graph over
$\partial B_1$:
\[
\partial\Otilde=\bigl\{(1+g(\theta))\theta:\theta\in\partial B_1\bigr\},
\]
with $\|g\|_{C^{1,\gamma}(\partial B_1)}\le M_1(N,R)$ and
$\|g\|_{L^\infty(\partial B_1)}\le C_H''(N,R)\,\alpha(\Otilde)^{1/(2N)}$.
The chain is asymmetry $\to$ Hausdorff (BDV Lemma~4.2 + density) $\to$
Alt--Caffarelli flatness (BDV Theorem~4.18) $\to$ tubular
patching. The threshold
$\alpha_{\rm graph}(N,R)=
\min\{\alpha_0,(h_F/C_H')^{2N},(h_F/C_{\rm Fr}')^2\}$ is explicit.
\end{theorem}

\subsection{Block III: Schauder bootstrap and interpolation}

\begin{theorem}[\textsc{SchauderInterpolation}; see
\texttt{SchauderInterpolation.tex}]\label{thm:I}
Fix $0<\gamma_0<\gamma\le1$. There exist
$M_5(N,R)$ and an explicit threshold
$\alpha_{\rm sph}(N,R,\gamma_0)$ (closed form in upstream constants)
such that the graph $g$ of Theorem~\ref{thm:G} satisfies
\[
\alpha(\Otilde)\le\alpha_{\rm sph}(N,R,\gamma_0)
\quad\Longrightarrow\quad
\|g\|_{C^{2,\gamma_0}(\partial B_1)}\le\delta_{\rm sph}(N,\gamma_0),
\]
where $\delta_{\rm sph}(N,\gamma_0)$ is the smallness threshold of
BDV Theorem~3.3. The method is the Kinderlehrer--Nirenberg hodograph
straightening of $\widetilde u$ on $\Otilde$, the oblique Schauder
bootstrap of Lieberman to $\|g\|_{C^{5,\gamma}}\le M_5(N,R)$, and the
Hölder interpolation $\|g\|_{C^{2,\gamma_0}}\le
C_{\rm GN}\|g\|_{L^\infty}^{1-\theta}\|g\|_{C^{5,\gamma}}^\theta$
with $\theta=(2+\gamma_0)/(5+\gamma)\le 1/2$.
\end{theorem}

\subsection{Block IV: nearly spherical closure}

\begin{theorem}[\textsc{NearlySphericalClosure}; see
\texttt{NearlySphericalClosure.tex}]\label{thm:N}
There is $c_*(N)>0$, given explicitly by
$c_*(N)=1/(32\,N^2\,C_3(N))$ where $C_3(N)$ is the BDV
Lemma~4.2(iii) constant, such that every nearly spherical
$\Otilde=(1+g)B_1$ with
\[
\|g\|_{C^{2,\gamma_0}}\le\delta_{\rm sph}(N,\gamma_0),\qquad
|\Otilde|=|B_1|,\qquad x_{\Otilde}=0,
\]
satisfies $Q_\alpha(\Otilde)\ge c_*(N)$, equivalently
$E(\Otilde)-E(B_1)\ge c_*(N)\alpha(\Otilde)$. Volume / barycenter
perturbations are absorbed by an $O(\|g\|_{L^2})$ correction at the
cost of a fixed dimensional factor. The proof is a direct substitution
into BDV Theorem~3.3
($E(\Om)-E(B_1)\ge(1/32N^2)\|g\|_{H^{1/2}}^2$) followed by the
trace-Poincar\'e comparison and BDV Lemma~4.2(iii)
($\alpha\le C_3\|g\|_{L^2}^2$).
\end{theorem}

\subsection{Block V: assembly to sharp Saint--Venant stability}

Composing Blocks I--IV gives the linear-in-$\alpha$ inequality
$E(\Om_0)-E(B_1)\ge q_*(N,R)\alpha(\Om_0)$ on a small-$\alpha$ regime,
with $q_*(N,R)=c_*(N)/(2C_{\rm sel}(N,R))$. Combining with BDV
Lemma~4.2(i) ($\Fr^2\le(C_1/\omega_N^2)\alpha$) gives the quadratic
form $E-E(B_1)\ge\kappa_{\rm near}(N,R)\Fr^2$. The far regime is
covered by BDV Lemma~5.2 ($E-E(B_1)\ge(C_8')^{-1}\Fr^4$), and the two
regimes glue at a threshold $a_{\rm near}=a_{\rm far}$ supplied by BDV
Lemma~4.2(ii).

\begin{theorem}[\textsc{SaintVenantAssembly}; see
\texttt{SaintVenantAssembly.tex}]\label{thm:SV}
There exists $c_{\rm SV}(N,R)>0$, given by~\eqref{eq:cSVdef} below,
such that for every open $\Om\subset B_R$ with $|\Om|=|B_1|$,
\begin{equation}\label{eq:SV}
   E(\Om)-E(B_1)\;\ge\;c_{\rm SV}(N,R)\,\Fr(\Om)^2.
\end{equation}
Concretely, with $\alpha_{\rm small}(N,R)
:=\tfrac12\min\{\eps_{\rm sel},\alpha_{\rm graph},\alpha_{\rm sph}\}$ and
$a_{\rm far}(N,R):=\alpha_{\rm small}(N,R)/(C_2(N,R)\omega_N)$,
\begin{equation}\label{eq:cSVdef}
   c_{\rm SV}(N,R)
   :=\min\!\Bigl\{\,
        \frac{q_*(N,R)\,\omega_N^2}{C_1(N)},\;\;
        \frac{a_{\rm far}(N,R)^2}{C_8'(N)}
   \Bigr\}.
\end{equation}
The exponent $2$ on $\Fr$ is sharp (ellipsoid example).
\end{theorem}

\subsection{Block VI: reduction to bounded sets}

The Saint--Venant assembly of Theorem~\ref{thm:SV} delivers a
constant $c_{\rm SV}(N,R)$ that depends on the confining radius
$R$. We remove this dependence by invoking BDV's constructive
truncation lemma~\cite[Lemma~5.3]{BDV}, which sits inside a chain of
three quantitative ingredients of \cite[\S5]{BDV} (all
constructive: De~Giorgi-type $L^{\infty}$ iteration, rearrangement,
and a cut-off competitor; no contradiction). Applied at the
dimensional radius $R_{*}(N):=\max\{d(N),2\}$ where $d(N)$ is the
diameter bound of BDV Lemma~5.3, the truncation step converts the
bounded-domain stability into a global one.

\begin{theorem}[\textsc{BoundedReduction}; see
\texttt{BoundedReduction.tex}]\label{thm:BR}
With $R_{*}(N):=\max\{d(N),2\}$ and
\begin{equation}\label{eq:cSVglob}
   c_{\rm SV}^{\rm glob}(N)\;:=\;
   \min\!\Bigl\{\,
      c_{\rm near}^{\rm glob}(N),\;
      \omega_{N}^{(N+2)/N}\,\delta(N)/16
   \Bigr\}>0,
\end{equation}
where
$c_{\rm near}^{\rm glob}(N)=\omega_{N}^{(N+2)/N}\bigl/\bigl[
4C_{9}(N)(1/\widetilde c_{\rm SV}(N)+C_{9}(N))\bigr]$
and $\widetilde c_{\rm SV}(N)=\omega_{N}^{-(N+2)/N}c_{\rm SV}(N,R_{*}(N))$,
for every open $\Om\subset\R^{N}$ with $|\Om|=|B_{1}|$,
\begin{equation}\label{eq:SVglob}
   E(\Om)-E(B_{1})\;\ge\;c_{\rm SV}^{\rm glob}(N)\,\Fr(\Om)^{2}.
\end{equation}
$c_{\rm SV}^{\rm glob}(N)$ depends only on $N$; the
$R$-dependence inherited from Theorem~\ref{thm:SV} is internalised at
the dimensional radius $R_{*}(N)$.
\end{theorem}

\subsection{Block VII: Kohler--Jobin transfer}

\begin{theorem}[\textsc{KohlerJobinTransfer}; see
\texttt{KohlerJobinTransfer.tex}]\label{thm:KJ}
Given the global sharp Saint--Venant inequality~\eqref{eq:SVglob},
for every open $\Om\subset\R^{N}$ of finite positive measure,
\begin{equation}\label{eq:FK}
   \Bigl(\frac{|\Om|}{|B_1|}\Bigr)^{2/N}
   \bigl(\lambda_1(\Om)-\lambda_1(B^*)\bigr)
   \;\ge\;c_{\rm FK}(N)\,\Fr(\Om)^2,
\end{equation}
where
\begin{equation}\label{eq:cFKdef}
   c_{\rm FK}(N)
   :=\min\!\Bigl\{\,
       \frac{4L_N\,c_{\rm SV}^{\rm glob}(N)}{(N+2)T_N},\;\;
       \frac{L_N(2^{2/(N+2)}-1)}{4}
   \Bigr\}.
\end{equation}
The proof uses the Kohler--Jobin inequality
$T(\Om)\lambda_1(\Om)^{(N+2)/2}\ge T(B^*)\lambda_1(B^*)^{(N+2)/2}$
(cited from \cite[Thm~3.3]{Brasco2014}), the pointwise transfer
$\lambda_1(\Om)\ge\lambda_1(B^*)(T(B^*)/T(\Om))^{2/(N+2)}$, the elementary
inequality $(1-s)^{-p}-1\ge p s$ for $p\in(0,1], s\in[0,1)$ (no
Taylor), and a two-case split on $T_N-T(\Om)$ versus $T_N/2$. Note
that the input~\eqref{eq:SVglob} is now unrestricted in $\Om$, so the
output~\eqref{eq:FK} is too.
\end{theorem}

\section{The assembled theorem}\label{sec:assembled}

Composing Theorems~\ref{thm:S}--\ref{thm:KJ}:

\begin{theorem}[Sharp quantitative Faber--Krahn stability, Plan~1 /
Route~A]\label{thm:final}
Let $N\ge2$. There exists a constant $c_{\rm FK}(N)>0$, depending
only on the dimension, given explicitly by~\eqref{eq:cFKdef} as a
composition of the named constants of
Theorems~\ref{thm:S}--\ref{thm:KJ}, such that for every open set
$\Om\subset\R^{N}$ of finite positive measure (where $B^*$ is the ball
of the same volume),
\[
   \boxed{\;
   \Bigl(\frac{|\Om|}{|B_1|}\Bigr)^{2/N}
   \bigl(\lambda_1(\Om)-\lambda_1(B^*)\bigr)
   \;\ge\;c_{\rm FK}(N)\,\Fr(\Om)^2.
   \;}
\]
Equivalently, $\Fr(\Om)\le C(N)\bigl((|\Om|/|B_1|)^{2/N}
(\lambda_1(\Om)-\lambda_1(B^*))\bigr)^{1/2}$,
and the exponent $1/2$ on the deficit is sharp (saturated by
volume-preserving ellipsoidal perturbations of $B^*$).
\end{theorem}

\begin{proof}
Combine Theorem~\ref{thm:SV} (sharp Saint--Venant stability in
$B_{R_{*}(N)}$) with Theorem~\ref{thm:BR} (\textsc{BoundedReduction},
removing the $R$-dependence) and Theorem~\ref{thm:KJ}
(Kohler--Jobin transfer). Theorem~\ref{thm:SV} is itself the result
of composing Theorems~\ref{thm:S}--\ref{thm:N} at the dimensional
radius $R=R_{*}(N)$ (the chain (S)$\to$(G)$\to$(I)$\to$(N) inside the
small-$\alpha$ regime, combined with BDV Lemma~5.2 in the far regime,
glued via the two directions of BDV Lemma~4.2).
\end{proof}

\section{What replaces the contradiction step}\label{sec:replaces}

BDV's original proof~\cite{BDV} establishes Theorem~\ref{thm:final}
via a sequential selection plus compactness contradiction: assume a
sequence $\Om_n$ violating the inequality with $\Fr(\Om_n)\to0$, run
the BDV penalised functional to select $\widetilde\Om_n$, extract an
$L^1$-convergent subsequence, and reach a contradiction with the
nearly spherical theorem at the limit. The compactness step makes the
final constant $c_{\rm FK}$ non-extractable.

Route~A replaces this with a \emph{single-set} version of the same
selection (Theorem~\ref{thm:S}), augmented by the quantitative graph
entry, Schauder bootstrap, and interpolation steps
(Theorems~\ref{thm:G}--\ref{thm:I}), so that the BDV nearly spherical
theorem (Theorem~\ref{thm:N}) can be applied to a single set
$\widetilde\Om$ rather than to a limit. The Saint--Venant assembly
(Theorem~\ref{thm:SV}) uses BDV's own Lemma~5.2 (proved by
rearrangement, not contradiction) for the far regime. The
\textsc{BoundedReduction} step (Theorem~\ref{thm:BR}) uses BDV's
constructive truncation lemma~\cite[Lemma~5.3]{BDV} to remove the
$R$-dependence (no contradiction there either). The Kohler--Jobin
transfer (Theorem~\ref{thm:KJ}) is unchanged from BDV.

Every constant in the chain depends only on $N$ and on the named
constants of BDV's quantitative ingredients (Lemmas~4.2, 4.5, 4.9,
4.12, 4.16, 5.1, 5.2, 5.3, Theorems~3.3, 4.18); none is set by a
compactness or contradiction argument.

\begin{thebibliography}{99}

\bibitem{BDV} L.~Brasco, G.~De~Philippis, B.~Velichkov,
\emph{Faber--Krahn inequalities in sharp quantitative form},
Duke Math.\ J.\ \textbf{164} (2015), 1777--1831 (arXiv:1306.0392).

\bibitem{FMP} N.~Fusco, F.~Maggi, A.~Pratelli,
\emph{The sharp quantitative isoperimetric inequality},
Ann.\ of Math.\ (2) \textbf{168} (2008), 941--980.

\bibitem{Brasco2014} L.~Brasco,
\emph{On torsional rigidity and principal frequencies: an invitation
to the Kohler--Jobin rearrangement technique},
ESAIM Control Optim.\ Calc.\ Var.\ \textbf{20} (2014), 315--338.

\bibitem{KN} D.~Kinderlehrer, L.~Nirenberg,
\emph{Regularity in free boundary problems},
Ann.\ Scuola Norm.\ Sup.\ Pisa Cl.\ Sci.\ (4) \textbf{4} (1977), 373--391.

\bibitem{Li86} G.~M.~Lieberman,
\emph{Mixed boundary value problems for elliptic and parabolic
differential equations of second order},
J.~Math.\ Anal.\ Appl.\ \textbf{113} (1986), 422--440.

\bibitem{Faber1923} G.~Faber,
\emph{Beweis, dass unter allen homogenen Membranen von gleicher
Fl\"ache und gleicher Spannung die kreisf\"ormige den tiefsten
Grundton gibt},
Sitzungsber.\ Bayer.\ Akad.\ Wiss.\ M\"unchen,
Math.-Phys.\ Kl.\ (1923), 169--172.

\bibitem{Krahn1925} E.~Krahn,
\emph{\"Uber eine von Rayleigh formulierte Minimaleigenschaft des
Kreises}, Math.\ Ann.\ \textbf{94} (1925), 97--100.

\bibitem{KawohlBook} B.~Kawohl,
\emph{Rearrangements and Convexity of Level Sets in PDE},
Lecture Notes in Math.\ \textbf{1150}, Springer, 1985.

\bibitem{HenrotBook} A.~Henrot,
\emph{Extremum problems for eigenvalues of elliptic operators},
Birkh\"auser, 2006.

\bibitem{BucurButtazzo} D.~Bucur, G.~Buttazzo,
\emph{Variational Methods in Shape Optimization Problems},
Progr.\ Nonlinear Differential Equations Appl.\ \textbf{65},
Birkh\"auser, 2005.

\bibitem{HansenNadirashvili1994} W.~Hansen, N.~Nadirashvili,
\emph{Isoperimetric inequalities in potential theory},
Potential Anal.\ \textbf{3} (1994), 1--14.

\bibitem{Melas1992} A.~D.~Melas,
\emph{The stability of some eigenvalue estimates},
J.\ Differential Geom.\ \textbf{36} (1992), 19--33.

\bibitem{Bhattacharya2001} T.~Bhattacharya,
\emph{Some observations about the first eigenvalue of the
$p$-Laplacian and its connections with asymmetry},
Electron.\ J.\ Differential Equations \textbf{35} (2001), 1--15.

\bibitem{FMPLaplace} N.~Fusco, F.~Maggi, A.~Pratelli,
\emph{Stability estimates for certain Faber--Krahn, isocapacitary and
Cheeger inequalities},
Ann.\ Sc.\ Norm.\ Super.\ Pisa Cl.\ Sci.\ (5) \textbf{8} (2009),
51--71.

\bibitem{BraDeP2017} L.~Brasco, G.~De~Philippis,
\emph{Spectral inequalities in quantitative form}, in:
\emph{Shape Optimization and Spectral Theory}, A.~Henrot (ed.),
De Gruyter Open, 2017, pp.~201--281.

\bibitem{CL2012} M.~Cicalese, G.~P.~Leonardi,
\emph{A selection principle for the sharp quantitative isoperimetric
inequality}, Arch.\ Ration.\ Mech.\ Anal.\ \textbf{206} (2012),
617--643.

\bibitem{FigalliMaggiPratelli} A.~Figalli, F.~Maggi, A.~Pratelli,
\emph{A mass transportation approach to quantitative isoperimetric
inequalities}, Invent.\ Math.\ \textbf{182} (2010), 167--211.

\bibitem{FuscoJulin} N.~Fusco, V.~Julin,
\emph{A strong form of the quantitative isoperimetric inequality},
Calc.\ Var.\ Partial Differential Equations \textbf{50} (2014),
925--937.

\bibitem{FuscoSurvey} N.~Fusco,
\emph{The quantitative isoperimetric inequality and related topics},
Bull.\ Math.\ Sci.\ \textbf{5} (2015), 517--607.

\bibitem{Maggi2008} F.~Maggi,
\emph{Some methods for studying stability in isoperimetric type
problems}, Bull.\ Amer.\ Math.\ Soc.\ \textbf{45} (2008), 367--408.

\bibitem{KohlerJobin} M.-T.~Kohler-Jobin,
\emph{Une m\'ethode de comparaison isop\'erim\'etrique de
fonctionnelles de domaines de la physique math\'ematique},
Z.\ Angew.\ Math.\ Phys.\ \textbf{29} (1978), 757--766.

\bibitem{AFM} E.~Acerbi, N.~Fusco, M.~Morini,
\emph{Minimality via second variation for a nonlocal isoperimetric
problem}, Comm.\ Math.\ Phys.\ \textbf{322} (2013), 515--557.

\bibitem{AC1981} H.~W.~Alt, L.~A.~Caffarelli,
\emph{Existence and regularity for a minimum problem with free
boundary}, J.\ Reine Angew.\ Math.\ \textbf{325} (1981), 105--144.

\bibitem{Fuglede} B.~Fuglede,
\emph{Stability in the isoperimetric problem for convex or nearly
spherical domains in $\mathbb R^n$},
Trans.\ Amer.\ Math.\ Soc.\ \textbf{314} (1989), 619--638.

\bibitem{AKN} M.~Allen, D.~Kriventsov, R.~Neumayer,
\emph{Sharp quantitative Faber--Krahn inequalities via the
Alt--Caffarelli--Friedman monotonicity formula} (preprint).

\bibitem{CFMP} A.~Cianchi, N.~Fusco, F.~Maggi, A.~Pratelli,
\emph{The sharp Sobolev inequality in quantitative form},
J.\ Eur.\ Math.\ Soc.\ \textbf{11} (2009), 1105--1139.

\bibitem{NeumayerWulff} R.~Neumayer,
\emph{A strong form of the quantitative Wulff inequality},
SIAM J.\ Math.\ Anal.\ \textbf{48} (2016), 1727--1772.

\bibitem{BrascoPratelli} L.~Brasco, A.~Pratelli,
\emph{Sharp stability of some spectral inequalities},
Geom.\ Funct.\ Anal.\ \textbf{22} (2012), 107--135.

\end{thebibliography}

\end{document}
